\section{Related Work}

Auxiliary contexts and demonstrations adapt a frozen predictor to a query
\citep{brown2020language,min2022rethinking}. Their quality and ordering matter,
and many methods retrieve or select examples before prediction
\citep{liu2022what,rubin2022learning,zhang2022active,su2023selective}. Other
work studies prompt-order sensitivity and calibration
\citep{lu2022fantastically,zhao2021calibrate}. Retrieval-augmented systems
condition predictions on external passages and records
\citep{guu2020realm,lewis2020retrieval,karpukhin2020dense}, and
retrieval scale and adaptation are studied in
\citep{borgeaud2022improving,izacard2023atlas,shi2024replug,asai2024selfrag}.
Recent work combines retrieval with adaptive chains and causal reasoning
\citep{acra2025retrieval,crgs2026causal}. Diversity-aware selection penalizes
repeated content or maximizes coverage \citep{carbonell1998use,gupta-etal-2023-coverage,levy2023diverse},
and listwise rerankers compare candidate sets directly
\citep{nogueira2019passage,ma2023zero,qin2024large,sun2023chatgpt,sachan2022improving}.
Long-context studies show that position and context composition change how a
model uses retrieved information \citep{liu2024lost}. These approaches
determine which information is relevant or diverse; they do not directly
estimate the change in end-task loss produced by adding a candidate to the
context already selected.

Outcome-aware selection moves closer to the final prediction objective by
valuing examples through their effect on model output. Hashimoto et al.\
estimate the incremental utility of a demonstration relative to an empty
prompt \citep{hashimoto-etal-2024-take}. Data valuation and influence
estimation assign credit to training examples through their effect on
predictions \citep{ghorbani2019data,koh2017understanding,pruthi2020estimating,ilyas2022datamodels}.
These methods characterize a candidate's standalone contribution. Our object
conditions on a nonempty, evolving selected set, so the value of the same
candidate can change after every routing decision.

Set-selection methods provide the structural counterpart to this decision
problem. Submodular objectives formalize coverage, diversity, and diminishing
returns \citep{lin2011class,nemhauser1978analysis,krause2014submodular}, while
determinantal point processes model diverse subsets probabilistically
\citep{kulesza2012determinantal}. Adaptive submodularity and influence
maximization extend these tools to sequential decisions
\citep{golovin2011adaptive,leskovec2007cost}, and streaming or contextual
variants address budgeted settings
\citep{badanidiyuru2014streaming,ross2013contextual}. Classical monotone
formulations reward additional elements; context utility in our setting can
also be negative because a candidate can increase prediction loss. Negative
transfer has been characterized in general transfer learning
\citep{pan2010survey,torrey2010transfer,wang2019characterizing} and mitigated
in sequential recommendation through multi-modal re-learning
\citep{ant2025negative}.

Routing and decision-focused learning treat selection as a learned decision
problem. Mixture-of-experts systems route inputs to expert subsets
\citep{jacobs1991adaptive,shazeer2017outrageously,zhou2022mixture}, and sparse
or soft routing architectures reduce computation while increasing capacity
\citep{fedus2022switch,lewis2021base,roller2021hash,puigcerver2024sparse}.
Routing networks select computation paths for multi-task learning
\citep{rosenbaum2018routing,ma2018modeling}, and recent systems apply dynamic
routing across language, vision-language-action, and collaborative-perception
tasks \citep{failuremode2026routing,dream2026routing,cogmoe2026sparse}.
Decision-focused learning connects prediction quality to downstream decision
quality \citep{bertsimas2020predictive,elmachtoub2022smart,wilder2019melding,mandi2024decision},
and task-based end-to-end learning follows the same principle
\citep{donti2017task}. Sequential selection is related to bandit exploration
and off-policy value estimation \citep{auer2002finite,dudik2011doubly}. R-MUR
applies this perspective to auxiliary contexts for a fixed predictor, where the
routing target is the marginal change in end-task loss and the predictor's
response becomes part of the available decision information.
